\fig{em2:lead-paths}{同一对端点的两条测量路径；两表均按从$a$到$b$计有向读数}{
\begin{tikzpicture}[scale=.85]
\fill[blue!6](-2.5,-1.3)rectangle(2.5,1.3);
\draw[world](-2.5,0)--(-2.5,1.3)--(-.35,1.3)
 (.35,1.3)--(2.5,1.3)--(2.5,0);
\draw[world](-2.5,0)--(-2.5,-1.3)--(-.35,-1.3)
 (.35,-1.3)--(2.5,-1.3)--(2.5,0);
\draw[world](0,1.3)circle(.35)node{$V_1$};
\draw[world](0,-1.3)circle(.35)node{$V_2$};
\fill(-2.5,0)circle(2pt)node[left]{$a$};
\fill(2.5,0)circle(2pt)node[right]{$b$};
\draw[flow,blue](-1.9,1.3)--(-1,1.3)node[midway,above]{$g_1$};
\draw[flow,red](-1.9,-1.3)--(-1,-1.3)node[midway,below]{$g_2$};
\node at(0,.1){$\otimes\ \bm n,\ \Phi_{12}$};
\node at(0,-2.05){$\Gamma=g_1+(-g_2)$：沿上路去、下路回};
\end{tikzpicture}}

图\ref{em2:lead-paths}中的两条路径都由$a$指向$b$，因此两支路的端电势项相同。设局部表体的感应可忽略，表内电阻降为$M_j=i_jR_{Vj}$，便有
\[
 M_1=V_a-V_b+\varepsilon_{g_1},\qquad
 M_2=V_a-V_b+\varepsilon_{g_2}.
\]
先作两式之差，再沿图中的闭合方向$\Gamma$使用法拉第定律，得到固定引线时
\[
 M_1-M_2=\int_{g_1}\bm E_f\cdot\dd\bm l
          -\int_{g_2}\bm E_f\cdot\dd\bm l
        =-\frac{\dd\Phi_{12}}{\dd t}.
\]
上路从左向右、下路反向返回，组成顺时针环路，故图中正法线指向纸里。读数差的符号须与这一法线一起确定。比较两次接线时还假定测量分流可忽略，或两次接线没有改变端点电势；若负载效应显著，必须分别重算被测网络。把两根引线贴近，正是缩小图中两路之间的面积。
